% !TEX program = lualatex
\documentclass[11pt,aspectratio=169]{beamer}
\usepackage[
  mode=beamer,
  theme=midnight,
  toc=true,
  toc-layout=page,
  numbering=none,
  density=compact
]{synaptic}

\SynapticCourseCode{MATH 201}
\SynapticCourse{Analysis}
\SynapticLectureNumber{1}
\SynapticTitle{Continuity, Compactness, and Finite Control}
\SynapticSubtitle{A visual route from local estimates to a global scale}
\SynapticShortTitle{Continuity and compactness}
\SynapticInstructor{Ada Lecturer}
\SynapticAffiliation{Department of Clear Mathematical Communication}
\SynapticSemester{Autumn 2026}
\SynapticDate{2 September 2026}
\SynapticTags{analysis \textbullet\ topology \textbullet\ proof design}
\SynapticSubject{A synaptic Beamer-mode design example}
\begin{synabstract}
  A restrained 16:9 presentation combining a clear title hierarchy, semantic
  cards, focused navigation, and mathematics designed for projection.
\end{synabstract}

\begin{document}
\SynapticMakeTitle

\section{From local to global}
\SynapticSectionFrame

\begin{frame}{One visual grammar}{Structure should survive projection and print}
  \begin{columns}[T,onlytextwidth]
    \begin{column}{0.48\textwidth}
      \begin{block}{Local information}
        For every point $x$, continuity supplies a neighbourhood and a scale
        $\delta_x$ adapted to that point.
      \end{block}
    \end{column}
    \begin{column}{0.48\textwidth}
      \begin{exampleblock}{Finite control}
        Compactness replaces the full family of neighbourhoods by finitely
        many representatives.
      \end{exampleblock}
    \end{column}
  \end{columns}

  \vspace{0.8ex}
  \begin{synlearninggoals}
    Distinguish local choices from a uniform estimate, then identify the exact
    step where compactness turns an open cover into finite data.
  \end{synlearninggoals}
\end{frame}

\begin{frame}{The compactness step}{A theorem card remains part of the argument}
  \begin{syntheorem}[Heine--Cantor]
    A continuous map $f\colon K\to Y$ from a compact metric space to a metric
    space is uniformly continuous.
  \end{syntheorem}

  \begin{synproof}
    Continuity gives neighbourhoods on which the oscillation of $f$ is small.
    Choose a finite subcover and take the minimum of its finitely many radii.
  \end{synproof}

  \[
    d_X(x,y)<\delta \quad\Longrightarrow\quad
    d_Y\bigl(f(x),f(y)\bigr)<\varepsilon .
  \]
\end{frame}

\section{A reusable proof pattern}
\SynapticSectionFrame

\begin{frame}{The proof as a four-step pipeline}
  \begin{enumerate}
    \item Start with the target tolerance $\varepsilon$.
    \item Build one local neighbourhood around every point.
    \item Extract a finite subcover using compactness.
    \item Convert the finite family into one global $\delta$.
  \end{enumerate}

  \begin{alertblock}{Common failure}
    Point-dependent choices $\delta_x$ do not by themselves establish uniform
    continuity; a finite reduction is the missing bridge.
  \end{alertblock}
\end{frame}

\begin{frame}{Takeaway}
  \begin{synlecturesummary}
    Continuity supplies local control. Compactness makes that control finite.
    A minimum over finitely many positive scales produces the uniform bound.
  \end{synlecturesummary}

  \vfill
  {\Large\bfseries\color{syn-primary}
    Local data $\longrightarrow$ finite data $\longrightarrow$ global scale.}
\end{frame}

\end{document}
